% © 2026 Ing. David Jaroš — CC BY-NC-ND 4.0
%
% This work is licensed under a Creative Commons Attribution-NonCommercial-NoDerivatives
% 4.0 International License (CC BY-NC-ND 4.0).
%
% License History: Earlier drafts (up to v0.3) were released under CC BY 4.0.
% From v0.4 onward, all material is released under CC BY-NC-ND 4.0 to protect
% the integrity of the theoretical work during ongoing academic development.
%
% See LICENSE.md for full license text.
%
% File: research_tracks/T3_ALPHA/integer_137_note.tex
% Purpose: Companion note — conditional integer-137 result, Gap G137-B formal
%          statement, and N_eff clarification.
% Status: Companion note only — NO new derivation steps.
%         Main result: [L1 conditional on B = B_phenom].
%         B_phenom NOT DERIVED. α NOT DERIVED.
% Date: 2026-06-11

\ifdefined\INCLUDEMODE
\else
\documentclass[12pt,a4paper]{article}
\usepackage[a4paper,margin=2.5cm]{geometry}
\usepackage{amsmath,amssymb,amsthm}
\usepackage{mathtools}
\usepackage{hyperref}
\usepackage{mdframed}
\usepackage{booktabs}
\usepackage{xcolor}
\usepackage{microtype}

\hypersetup{
  colorlinks=true,
  linkcolor=blue!60!black,
  citecolor=green!50!black,
  urlcolor=blue!60!black
}

\newtheorem{theorem}{Theorem}[section]
\newtheorem{lemma}[theorem]{Lemma}
\newtheorem{proposition}[theorem]{Proposition}
\newtheorem{corollary}[theorem]{Corollary}
\theoremstyle{definition}
\newtheorem{definition}[theorem]{Definition}
\theoremstyle{remark}
\newtheorem{remark}[theorem]{Remark}

\newmdenv[backgroundcolor=yellow!20, linecolor=orange, linewidth=1.5pt]{gapbox}
\newmdenv[backgroundcolor=green!15, linecolor=green!60!black, linewidth=1.5pt]{resultbox}
\newmdenv[backgroundcolor=red!10, linecolor=red!60, linewidth=1.5pt]{deadendbox}
\newmdenv[backgroundcolor=blue!8, linecolor=blue!50, linewidth=1.5pt]{insightbox}

\newcommand{\PROVED}{\textbf{[Proved]}}
\newcommand{\OPEN}{\textbf{[Open]}}
\newcommand{\Lzero}{\textbf{[L0]}}
\newcommand{\Lone}{\textbf{[L1]}}
\newcommand{\Ltwo}{\textbf{[L2]}}
\newcommand{\MC}{\textbf{[MC]}}
\newcommand{\NUM}{\textbf{[NUM]}}
\newcommand{\OBS}{\textbf{[OBS]}}
\newcommand{\COND}{\textbf{[L1 conditional]}}
\newcommand{\NODERIVED}{\textbf{NOT DERIVED}}

\newcommand{\CC}{\mathbb{C}}
\newcommand{\HH}{\mathbb{H}}
\newcommand{\RR}{\mathbb{R}}

\title{\textbf{Conditional Integer-137 Result in UBT}\\[6pt]
  \large Structural evidence, open gap G137-B, and N\textsubscript{eff}
  clarification}
\author{Ing.\ David Jaroš}
\date{2026-06-11}

\begin{document}
\maketitle

\begin{abstract}
We record the conditional result that the UBT winding-mode effective potential
$V_{\mathrm{eff}}(n) = n^2 - B\cdot n\ln n$ has its unique integer minimum at
$n^* = 137$ when $B = B_{\mathrm{phenom}} \approx 46.298$.  This result is
proved at level \COND\ conditional on $B$.  The coefficient $B_{\mathrm{phenom}}$
is \textbf{not derived} from the UBT action $S[\Theta]$; deriving it is Gap G137-B,
which remains open after six NO-GO routes.  The T3\_ALPHA track has been downgraded
from CONDITIONAL to STRUCTURAL EVIDENCE.  We state the three sub-gaps (G137-B-i,
ii, iii) precisely, record the six NO-GO route residuals, and clarify which
$N_{\mathrm{eff}}$ enters the integer-137 claim.
\textbf{Alpha is not derived.}
\end{abstract}

\tableofcontents

% ==============================================================================
\section{Main Result (Conditional)}
\label{sec:main}
% ==============================================================================

\begin{resultbox}
\begin{theorem}[Integer-137 from UBT effective potential, \COND]
\label{thm:integer_137}
Define the UBT winding-mode effective potential
\begin{equation}
\label{eq:Veff}
  V_{\mathrm{eff}}(n) = n^2 - B\cdot n\ln n, \qquad n \in \mathbb{Z}_{>0}.
\end{equation}
For $B = B_{\mathrm{phenom}} \approx 46.298$, the function $V_{\mathrm{eff}}$
has a unique integer minimum at
\begin{equation}
  n^* = \arg\min_{n \in \mathbb{Z}_{>0}} V_{\mathrm{eff}}(n) = 137.
\end{equation}
This is a prime number.  The result is at level \Lone\ conditional on
$B = B_{\mathrm{phenom}}$.
\end{theorem}
\end{resultbox}

\begin{proof}[Proof sketch]
The continuous minimum of $V_{\mathrm{eff}}(x) = x^2 - Bx\ln x$ is obtained by
setting $\dd V_{\mathrm{eff}}/\dd x = 0$:
\[
  2x - B(\ln x + 1) = 0,
\]
giving $x^* \approx e^{B/2 - 1/2}$ for large $B$.  For $B = 46.298$, this
gives $x^* \approx 136.8$, whose nearest integer is $n^* = 137$.  Checking
$V_{\mathrm{eff}}(136) < V_{\mathrm{eff}}(137)$ vs $V_{\mathrm{eff}}(138)$
confirms 137 is the unique integer attractor.  The primality of 137 provides
additional stability (prime attractor argument).  Full derivation chain:
\texttt{canonical/alpha/ALPHA\_MASTER\_STATUS.md} and
\texttt{canonical/alpha/alpha\_best\_route.tex}.
\end{proof}

\begin{gapbox}
\textbf{Conditionality}: The coefficient $B = B_{\mathrm{phenom}} \approx 46.298$
is \textbf{not derived} from $S[\Theta]$.  It is an observed numerical value
that gives the correct attractor.  Deriving $B_{\mathrm{phenom}}$ from first
principles is Gap G137-B (Section~\ref{sec:gap}).  Until G137-B is closed,
Theorem~\ref{thm:integer_137} is \textbf{conditional evidence only}, not a
derivation of the fine structure constant.
\end{gapbox}

% ==============================================================================
\section{The B Coefficient: What Is and Is Not Proved}
\label{sec:B}
% ==============================================================================

The coefficient $B$ in~\eqref{eq:Veff} enters through three distinct routes.
Their status is:

\begin{center}
\begin{tabular}{lll}
\toprule
\textbf{Claim} & \textbf{Level} & \textbf{Source} \\
\midrule
$B = 12^{3/2}\cdot(2\eta(i))^{1/4}$ (algebraic identity) &
  \Lzero & \texttt{canonical/alpha/ALPHA\_MASTER\_STATUS.md} \\
Numerical match $B_{\mathrm{Ram}} \approx 46.28 \approx B_{\mathrm{phenom}}$ &
  \OBS & \texttt{canonical/alpha/ALPHA\_MASTER\_STATUS.md} \\
Algebraic identity $\eta^{-2}\theta_3\theta_4^2 = 2\eta(i)$ &
  \Lzero & v42 result \\
$N_{\mathrm{eff}}^{\mathrm{twist}} = 12$ (SU(2)-twist route) &
  \Lone & \texttt{canonical/n\_eff/step2\_AUDIT.tex} \\
$B_{\mathrm{phenom}}$ from $S[\Theta]$ without $\alpha$ input &
  \textbf{NOT DERIVED} & Gap G137-B open \\
\bottomrule
\end{tabular}
\end{center}

\begin{insightbox}
\textbf{Status of $B$}: The Ramanujan form $B_{\mathrm{Ram}} = 12^{3/2}(2\eta(i))^{1/4}$
is an \OBS\ algebraic coincidence: $B_{\mathrm{Ram}} \approx 46.281$ vs
$B_{\mathrm{phenom}} \approx 46.298$, a 0.037\% discrepancy.  Whether this
coincidence has a structural explanation from $S[\Theta]$ is precisely
what Gap G137-B asks.  No claim is made that $B$ is derived.
\end{insightbox}

% ==============================================================================
\section{Gap G137-B: Formal Statement}
\label{sec:gap}
% ==============================================================================

Gap G137-B is the central open problem of the T3\_ALPHA track.  The following
three sub-gaps are copied directly from
\texttt{research\_tracks/T3\_ALPHA/mellin\_insertion\_B.tex} §formal-gap.

\begin{gapbox}
\textbf{Gap G137-B: Derive $B_{\mathrm{phenom}} \approx 46.298$ from $S[\Theta]$
without using $\alpha$ as input.}\\[6pt]
Three sub-gaps obstruct a complete derivation:
\begin{itemize}
\item \textbf{G137-B-i} [OPEN/MC]: Volumetric factorization
  $W_{\mathrm{eff}} = N_{\mathrm{eff}}^{3/2}\cdot f(Z_b)$ — the Mellin
  insertion structure requires a factorization of the partition function
  that has not been proved.
\item \textbf{G137-B-ii} [OPEN/MC]: $Z_{1\mathrm{real}} = 2\eta(i)$ from the
  NS-sector heat kernel — identifying the single-real-boson partition function
  with the Dedekind $\eta$ at the self-dual point has not been derived from the
  UBT action.
\item \textbf{G137-B-iii} [OPEN/MC]: $N_{\mathrm{eff}}^{1/2}$ from $T^3$ volume
  at the self-dual point — the volume factor that converts the zero-mode
  normalisation to $B_{\mathrm{phenom}}$ has not been determined from $S[\Theta]$.
\end{itemize}
All three sub-gaps must be closed simultaneously to derive $B_{\mathrm{phenom}}$
from first principles.
\end{gapbox}

\subsection{Six NO-GO routes}

The table below records the six routes exhausted in the time-box for Gap G137-B,
with their numerical residuals.  These are taken from
\texttt{research\_tracks/T3\_ALPHA/mellin\_insertion\_B.tex}.

\begin{center}
\begin{tabular}{llll}
\toprule
\textbf{\#} & \textbf{Route} & \textbf{Verdict} & \textbf{Residual / ratio} \\
\midrule
1 & T-duality & [NO-GO] & Algebraic obstruction; $B \neq B_{\mathrm{phenom}}$ \\
2 & Orbifold $\mathbb{Z}_{\mathrm{orb}}(i)$ & [NO-GO] & $\approx 1.302 \neq \sqrt{2}$ \\
3 & Direct $M^4$ functional determinant ($\zeta$-reg.) & [NO-GO] & Residual $\sim 1.14\times10^4$ \\
4 & Casimir $T^3$ & [NO-GO] & Residual $\sim 9.0\times10^3$ \\
5 & Eisenstein $E_4(i)$ & [NO-GO] & Ratio $B_{E_4}/B_{\mathrm{phenom}} \approx 0.920$ \\
6 & Mock theta $f(e^{-\pi})$ & [NO-GO] & Ratio $\approx 0.907$ \\
\bottomrule
\end{tabular}
\end{center}

\begin{deadendbox}
\textbf{Time-box expired.}  All six routes are NO-GO.  No further routes for
Gap G137-B will be attempted without explicit instruction.  The A\_PRIME modular
bootstrap route (A\_PRIME) has been exhausted.  T3\_ALPHA has been formally
downgraded to STRUCTURAL EVIDENCE.  See
\texttt{research\_tracks/T3\_ALPHA/mellin\_insertion\_B.tex} for full details.
\end{deadendbox}

% ==============================================================================
\section{N\textsubscript{eff} Clarification}
\label{sec:neff}
% ==============================================================================

Two independent routes to $N_{\mathrm{eff}}$ are tracked in
\texttt{canonical/n\_eff/step2\_AUDIT.tex}:

\begin{center}
\begin{tabular}{llll}
\toprule
\textbf{Route} & \textbf{Value} & \textbf{Level} & \textbf{Source} \\
\midrule
SU(2)-twist & $N_{\mathrm{eff}}^{\mathrm{twist}} = 12$ & \Lone &
  \texttt{step2\_AUDIT.tex} \\
Scalar loop & $N_{\mathrm{eff}}^{\mathrm{loop}} = 3$ & \Lone &
  \texttt{step2\_AUDIT.tex} \\
Identification twist = loop & \textbf{[OPEN/MC frozen]} &
  — & \texttt{step2\_AUDIT.tex §rem:neff\_alpha\_dependency} \\
\bottomrule
\end{tabular}
\end{center}

\begin{resultbox}
\textbf{Which $N_{\mathrm{eff}}$ enters the integer-137 result?}

The conditional result of Theorem~\ref{thm:integer_137} uses
$N_{\mathrm{eff}}^{\mathrm{twist}} = 12$ [L1] via the A\_PRIME route.

The scalar-loop route $N_{\mathrm{eff}}^{\mathrm{loop}} = 3$ [L1] gives
$B_0^{\mathrm{loop}} = 2\pi$, yielding $n^* \approx 10.5 \neq 137$.

Therefore:
\begin{enumerate}
\item The integer-137 companion note is conditional on the
  \textbf{SU(2)-twist route specifically}.
\item The identification $N_{\mathrm{eff}}^{\mathrm{twist}} = N_{\mathrm{eff}}^{\mathrm{loop}}$
  is \textbf{NOT} required for the integer-137 conditional result.
\item The identification would only be required for a first-principles derivation
  of $B_0$ from the scalar effective action — which remains part of Gap G137-B.
\end{enumerate}
Status of identification: \textbf{[OPEN/MC frozen]}.  See
\texttt{canonical/n\_eff/step2\_AUDIT.tex §rem:neff\_alpha\_dependency}.
\end{resultbox}

% ==============================================================================
\section{Structural Significance}
\label{sec:structure}
% ==============================================================================

\begin{insightbox}
\textbf{Why the integer-137 result is structurally significant, even as
structural evidence only.}

\begin{enumerate}
\item \textbf{Prime attractor}: $n^* = 137$ is prime.  The effective potential
  $V_{\mathrm{eff}}(n)$ has enhanced stability at prime $n^*$, because
  the winding-mode spectrum has no sub-harmonic resonance at a prime.

\item \textbf{Modular corroboration}: The modular group index
  $\mu(\Gamma_0(137))/3 \approx 46.00$ is within 0.64\% of $B_{\mathrm{phenom}}$,
  suggesting a structural connection between the prime attractor and the
  congruence subgroup $\Gamma_0(137)$.  This is an \OBS\ observation, not a
  derivation.

\item \textbf{No fine-tuning}: The potential $V_{\mathrm{eff}}(n) = n^2 - Bn\ln n$
  has no free parameters once $B$ is fixed.  Its minimum at $n^* = 137$ is a
  \emph{discrete} attractor — nearby values $n = 136, 138$ give strictly larger
  $V_{\mathrm{eff}}$.

\item \textbf{$\alpha$ is not derived}: The fine structure constant
  $\alpha \approx 1/137.036$ is \textbf{not} reproduced by this result.  Only
  the integer part $\alpha^{-1}_{\mathrm{bare}} = 137$ is structurally motivated.
  The renormalisation-group running correction $0.036$ is not addressed.
\end{enumerate}

\textbf{T3\_ALPHA track status}: STRUCTURAL EVIDENCE (downgraded from
CONDITIONAL on 2026-06-11 after six NO-GOs on Gap G137-B).  Publication target:
short companion note recording the conditional result and gap structure.
\end{insightbox}

% ==============================================================================
\section{Summary of T3\_ALPHA Status}
\label{sec:summary}
% ==============================================================================

\begin{center}
\begin{tabular}{lll}
\toprule
\textbf{Claim} & \textbf{Level} & \textbf{Status} \\
\midrule
$B = 12^{3/2}(2\eta(i))^{1/4}$ identity & \Lzero & PROVED \\
$n^*(B_{\mathrm{phenom}}) = 137$ & \COND & PROVED conditional on $B$ \\
$N_{\mathrm{eff}}^{\mathrm{twist}} = 12$ & \Lone & PROVED \\
$N_{\mathrm{eff}}^{\mathrm{loop}} = 3$ (independent) & \Lone & PROVED \\
Integer-137 uses $N_{\mathrm{eff}}^{\mathrm{twist}} = 12$,
  not $N_{\mathrm{eff}}^{\mathrm{loop}} = 3$ & \Lone & PROVED \\
$N_{\mathrm{eff}}^{\mathrm{twist}} = N_{\mathrm{eff}}^{\mathrm{loop}}$ & \textbf{[OPEN/MC frozen]} & NOT PROVED \\
G137-B-i, ii, iii (three sub-gaps) & \textbf{[OPEN/MC]} & NOT PROVED \\
$B_{\mathrm{phenom}}$ from $S[\Theta]$ & \textbf{NOT DERIVED} & — \\
$\alpha \approx 1/137.036$ from UBT & \textbf{NOT DERIVED} & — \\
\bottomrule
\end{tabular}
\end{center}

\begin{deadendbox}
\textbf{T3\_ALPHA: STRUCTURAL EVIDENCE}\\[4pt]
$\alpha$ is NOT derived in UBT as of 2026-06-11.  The integer-137 result is
conditional evidence only.  No new derivation routes for $\alpha$ will be
opened without explicit instruction.
\end{deadendbox}

% ==============================================================================
\begin{thebibliography}{99}

\bibitem{jaros2026_gr}
  Ing.~D.~Jaroš,
  \textit{General Relativity from Unified Biquaternion Theory},
  \texttt{papers/UBT\_GR\_Submission.tex} (2026).

\bibitem{jaros2026_gauge}
  Ing.~D.~Jaroš,
  \textit{Standard Model Gauge Structure from Unified Biquaternion Theory},
  \texttt{papers/UBT\_Gauge\_Submission.tex} (2026).

\bibitem{alpha_master}
  Ing.~D.~Jaroš,
  \textit{Alpha Master Status},
  \texttt{canonical/alpha/ALPHA\_MASTER\_STATUS.md} (2026).

\bibitem{mellin_nogo}
  Ing.~D.~Jaroš,
  \textit{Gap G137-B: Mellin Insertion — Formal No-Go Record},
  \texttt{research\_tracks/T3\_ALPHA/mellin\_insertion\_B.tex} (2026).

\bibitem{neff_audit}
  Ing.~D.~Jaroš,
  \textit{N\_eff Audit (Step 2)},
  \texttt{canonical/n\_eff/step2\_AUDIT.tex} (2026).

\end{thebibliography}

\ifdefined\INCLUDEMODE
\else
\end{document}
\fi
